\chapter{Reliability}
\label{ch:reliability}

A phone is an unreliable place to run a real-time app. The network drops and returns, the operating
system suspends apps in the background, Realtime connections break and rejoin, and the user may close the
app at any moment. This chapter describes how Yap stays correct under those conditions. The guiding rule
is that the server is the source of truth: the client may fall behind, but it always has a way to ask the
server where things stand and catch up.

\section{Failure model}

The design assumes five kinds of failure:

\begin{itemize}
  \item \textbf{Lost Realtime events.} Realtime does not replay changes sent while a channel was
    disconnected, so any event can be missed.
  \item \textbf{Background suspension.} Without push notifications, an app in the background receives
    nothing.
  \item \textbf{Network loss.} Any call can fail, including a call whose effect matters to someone else,
    such as going offline.
  \item \textbf{Restart.} The app can be closed while live, waiting for an answer or in a chat.
  \item \textbf{Races.} Two users can act on the same request at the same moment, and slow answers can
    arrive after newer ones.
\end{itemize}

\section{Realtime subscriptions}
\label{sec:realtime}

The client opens at most three Realtime channels of its own and one per waiting request
(\cref{fig:realtime}). Each is created by the API layer and subscribed through \code{subscribeAsMe}, which
first passes the user's current token to Realtime. Without that step, a session created by some sign-in
paths (the password-reset code, for one) would join channels without the user's token; the database would
then refuse the filter and nothing would arrive. Every subscription accepts an \code{onSubscribed}
callback that runs on the first join and again after every reconnect, which is the moment to fetch
whatever may have been missed.

\begin{figure}[ht]
  \centering
  % Realtime channel map: from table changes to store state (Figure: realtime map).
\begin{tikzpicture}[
    tbl/.style={box, serverbox, font=\ttfamily\scriptsize, text width=22mm, align=center},
    ch/.style={box, accentbox, font=\sffamily\scriptsize, text width=35mm, align=flush left, minimum height=10mm},
    hd/.style={box, clientbox, font=\sffamily\scriptsize, text width=26mm, align=flush left, minimum height=10mm},
    st/.style={box, font=\sffamily\scriptsize, text width=22mm, align=flush left, minimum height=10mm},
    lab/.style={font=\sffamily\tiny, text=muted, fill=white, inner sep=1pt},
  ]
  \foreach \x/\t in {1.3/TABLE, 6.6/CHANNEL AND FILTER, 10.75/HANDLER IN THE STORE, 14.55/STATE} {
    \node[layerlabel, anchor=south, inner sep=0pt] at (\x,0.75) {\t};}

  \node[tbl, minimum height=30mm] (req) at (1.3,-1.0) {chat\_requests};
  \node[tbl, minimum height=19mm] (msg) at (1.3,-4.15) {messages};
  \node[font=\sffamily\tiny, text=muted, align=center, text width=26mm, anchor=north] at (msg.south) {both in \texttt{supabase\_realtime}, replica identity full};

  \begin{scope}[on background layer]
    \fill[accentsoft!45] (3.25,0.6) rectangle (4.05,-5.4);
  \end{scope}
  \node[font=\sffamily\tiny, text=accentdark, rotate=90, align=center] at (3.65,-2.4) {select policies: only rows the subscriber may read};

  \node[ch] (c1) at (6.6,0) {\texttt{incoming-requests}\\INSERT, \texttt{to\_user = me}};
  \node[ch] (c2) at (6.6,-1.5) {\texttt{request-<id>}\\UPDATE, \texttt{id = <id>}};
  \node[ch, minimum height=14mm] (c3) at (6.6,-3.15) {\texttt{chat-<id>}\\INSERT, \texttt{request\_id = <id>}\\UPDATE request, \texttt{id = <id>}};
  \node[ch] (c4) at (6.6,-4.8) {\texttt{messages-event\_id-<id>}\\INSERT, \texttt{event\_id = <id>}};

  \node[hd] (h1) at (10.75,0) {\texttt{onIncoming}: fetch the sender's card};
  \node[hd] (h2) at (10.75,-1.5) {\texttt{waitForAnswer}: accepted, or not this time};
  \node[hd, minimum height=14mm] (h3) at (10.75,-3.15) {\texttt{addChatMessages}, merged by ID; \texttt{markChatEnded}};
  \node[hd] (h4) at (10.75,-4.8) {\texttt{updateEvent}: merge into the event's messages};

  \node[st] (s1) at (14.55,0) {\texttt{incoming}; \texttt{/incoming} opens};
  \node[st] (s2) at (14.55,-1.5) {\texttt{outgoing}; chat opens};
  \node[st, minimum height=14mm] (s3) at (14.55,-3.15) {\texttt{chat.messages}, \texttt{chat.status}};
  \node[st] (s4) at (14.55,-4.8) {\texttt{events[i]} \texttt{.messages}};

  \draw[push] (req.east |- c1) -- (c1);
  \draw[push] (req.east |- c2) -- (c2);
  \draw[push] (req.east |- 0,-2.75) -- (c3.west |- 0,-2.75);
  \draw[push] (msg.east |- 0,-3.55) -- (c3.west |- 0,-3.55);
  \draw[push] (msg.east |- c4) -- (c4);
  \foreach \i in {1,2,3,4} {\draw[flow] (c\i) -- (h\i); \draw[flow] (h\i) -- (s\i);}
\end{tikzpicture}

  \caption{Realtime channel map. Two published tables feed four channels; select policies decide which rows
    each subscriber receives; each channel's handler updates one part of the store.
    \Cref{tab:subscriptions} gives when each channel opens and closes.}
  \label{fig:realtime}
\end{figure}

The store keeps exactly one of each of its channels and closes all of them on sign-out. Opening a new
one-to-one chat or group chat closes the previous one first, so a channel is never leaked.

\section{Catch-up}
\label{sec:catch-up}

Because events can be missed, the store has a single procedure, \code{catchUp}, that asks the server
directly where things stand. It runs when the app returns to the foreground, when the network returns,
when the incoming-requests channel re-joins after a disconnection, and when Auth refreshes the token after
a failed load (\cref{fig:catch-up}). Calls never overlap: a call that arrives while one is running makes
it run once more when it finishes.

\begin{figure}[ht]
  \centering
  % The catch-up that runs when the app returns, reconnects or a channel re-joins (Figure: catch-up).
\begin{tikzpicture}[
    node distance=5mm and 8mm,
    st/.style={step=52mm, font=\sffamily\scriptsize},
    trig/.style={box, clientbox, font=\sffamily\scriptsize, text width=27mm, align=center},
    yes/.style={edgelabel, fill=none, font=\sffamily\scriptsize\itshape, text=muted},
    outcome/.style={terminal, font=\sffamily\scriptsize, text width=22mm, align=center},
  ]
  \node[trig] (t1) at (0,0) {app returns to\\the foreground};
  \node[trig] (t2) at (3.3,0) {network comes\\back};
  \node[trig] (t3) at (6.6,0) {incoming channel\\re-joins};
  \node[trig] (t4) at (9.9,0) {token refreshed\\after a failed load};

  \node[decision, aspect=2.6] (running) at (4.95,-1.75) {catch-up\\running?};
  \node[outcome, right=12mm of running] (again) {run once more\\when it ends};
  \node[decision, aspect=2.6, below=of running] (failed) {profile\\failed to load?};
  \node[outcome, right=12mm of failed] (reload) {retry the full\\load, then stop};
  \node[decision, aspect=2.6, below=of failed] (profile) {profile\\loaded?};
  \node[outcome, right=12mm of profile] (stop) {stop: signed out\\or still setting up};
  \node[st, below=of profile] (offline) {retry a pending ``go offline''};
  \node[decision, aspect=2.6, below=of offline] (extras) {block list and\\status loaded?};
  \node[st, text width=34mm, right=12mm of extras] (load) {load the block list and restore from \texttt{my\_status}; stop if still offline};
  \node[st, below=of extras] (status) {\texttt{my\_status}: drop a live session the server has ended; open a chat the server has but this phone does not};
  \node[st, below=of status] (chatcheck) {request status of the open chat: mark it ended unless still \texttt{accepted}};
  \node[st, below=of chatcheck] (incoming) {look for a waiting request addressed to the user, unless one is on screen or a chat is open};
  \node[st, below=of incoming] (sync) {reload the open chat's newest 200 messages};

  \foreach \t in {t1, t2, t3, t4} {\draw[flow] (\t.south) -- (running);}
  \draw[flow] (running) -- node[yes, above] {yes} (again);
  \draw[flow] (running) -- node[yes, right] {no} (failed);
  \draw[flow] (failed) -- node[yes, above] {yes} (reload);
  \draw[flow] (failed) -- node[yes, right] {no} (profile);
  \draw[flow] (profile) -- node[yes, above] {no} (stop);
  \draw[flow] (profile) -- node[yes, right] {yes} (offline);
  \draw[flow] (offline) -- (extras);
  \draw[flow] (extras) -- node[yes, right] {yes} (status);
  \draw[flow] (extras) -- node[yes, above] {no} (load);
  \draw[flow] (status) -- (chatcheck);
  \draw[flow] (load.south) |- (chatcheck.east);
  \draw[flow] (chatcheck) -- (incoming);
  \draw[flow] (incoming) -- (sync);
\end{tikzpicture}

  \caption{The catch-up. Each step tolerates failure; whatever fails is tried again by the next
    catch-up.}
  \label{fig:catch-up}
\end{figure}

Each step repairs one kind of drift between the phone and the server:

\begin{itemize}
  \item A pending request to go offline is sent again (\cref{sec:offline}).
  \item If the profile has not loaded, because the user is signed out or still setting up, the catch-up
    stops here. If the block list and status never loaded, they are loaded now instead of the next step.
  \item \code{my_status} reveals a live session the server has ended while the app was asleep; the app
    then shows \ui{Your time's up}. It also reveals a chat that was accepted just before the connection
    dropped.
  \item A query of the open chat's request status (\code{api.requestStatus}) reveals a chat the other
    user ended or deleted with their account: anything other than \code{accepted} marks it ended.
  \item A query for the newest waiting request addressed to the user finds a request that arrived while
    the app was in the background.
  \item The open chat's history is reloaded and merged by message ID.
\end{itemize}

The catch-up never undoes something the user did while it was running. It drops the live session only if
the session on the phone is the same object it was before the call, so a fresh \ui{Go live} is not
cancelled by a stale answer, and it opens a chat only if no chat changed on the phone in the meantime.

\section{Offline behaviour}
\label{sec:offline}

\Cref{tab:offline} lists how each action behaves without a connection. A banner tells the user the phone
is offline whenever the operating system reports no connection or no internet.

\begingroup\tablefont
\begin{longtable}{@{}L{46mm}L{\dimexpr\linewidth-46mm-2\tabcolsep\relax}@{}}
  \caption{Behaviour without a connection.}
  \label{tab:offline}\\
  \toprule
  \textsf{\textbf{Situation}} & \textsf{\textbf{Behaviour}} \\
  \midrule
  \endfirsthead
  \toprule
  \textsf{\textbf{Situation}} & \textsf{\textbf{Behaviour}} \\
  \midrule
  \endhead
  \bottomrule
  \endfoot
  Launch with a stored session & The profile cannot load, so \ui{Can't reach Yap} replaces the app, with \ui{Try again}; the load is retried when the token refreshes or a catch-up runs \\
  Go live, say hi, accept, send a message, create, join or delete an event & The call fails and an alert explains; nothing changes on the server. A failed \ui{Accept} also closes the request screen, and the request later expires \\
  Go offline by tapping the red button & The phone shows the user offline at once; \ui{Almost offline} appears, and the call is retried at the next catch-up \\
  Live session reaches its end time & The phone shows the user offline; the server's minute job removes the session on time anyway \\
  Nearby refresh & The last list stays on screen \\
  Waiting for an answer & Realtime and the poll fail quietly; the timer shows \ui{Not this time} after the end time; if the request was in fact accepted, the catch-up opens the chat \\
  Chat accepted but not yet opened & An alert says the chat opens as soon as the phone is back online; the catch-up opens it \\
  Messages, requests or a chat ending missed while away & Recovered by the channel re-joining and by the catch-up \\
  Map & The existing pins stay \\
  A person's further photos & The profile shows without them \\
  Sign out & Refused with an explanation, because ending the chat and going offline must reach the server first \\
\end{longtable}
\endgroup

\section{Ordering and races}
\label{sec:ordering}

Several guards keep late or duplicate work from corrupting state.

\paragraph{On the server} \code{respond}, \code{cancel_request} and \code{end_chat} lock the request row
before deciding, so concurrent actions on one request apply in sequence. A partial unique index makes a
second waiting request from one sender impossible even if two calls race. \code{respond} withdraws the
recipient's own waiting request before checking for an existing chat, so an answer landing at the same
moment is seen. A cancel that arrives just after an accept ends the new chat instead of leaving one user
in a chat the other never sees.

\paragraph{Double taps} Every button, icon button, row and card ignores a second tap while the action
the first one started is still running (the button shows its spinner meanwhile), and for half a second
after a tap that only opens or closes a screen (\code{useGuardedPress} in
\code{src/components/ui/press-feedback.ts}), so a double tap cannot submit twice, open a screen twice or go
back two screens. Forms that a keyboard key or a completed code can also submit run one attempt at a time
(\code{useOneAtATime}). The server makes the repeats that still get through harmless: saying hi again to
the person already waited on returns the same request, a repeated report keeps the first, and joining,
blocking and going live again change nothing.

\paragraph{In the store}
\begin{itemize}
  \item A say-hi on its way to the server blocks a second one, and a sign-out or account deletion
    already under way makes a second tap of the same action wait for it. The other action is refused
    with an error rather than waiting, so a sign-out never passes for a deletion.
  \item Sign-out resets every counter and in-flight marker below, so an answer meant for the previous
    account lands nowhere.
  \item Counters on nearby refreshes, event loads and group chat visits discard any answer that arrives
    after a newer request has started or the screen has closed.
  \item \code{handledRequestIds} and \code{closedChatIds} stop a catch-up, which may have read the server a
    moment earlier, from reviving a request already answered or a chat already closed on this phone.
  \item \code{openingChatId} stops the same chat being opened twice when Realtime and a catch-up report it
    together.
  \item Chats subscribe before loading history and merge messages by ID, so no message is lost between
    the two steps and none is shown twice.
  \item Going live waits for an in-flight retry of going offline, so the retry cannot delete the new
    session.
  \item \code{init()} runs once per process, however often the root layout mounts.
\end{itemize}

\section{Time}
\label{sec:time}

The server owns every deadline. \code{go_green} returns the session's start and end as the server set
them, and \code{send_request} returns the request's end time; the countdowns on screen count towards those
instants. A phone's clock can be wrong, set by hand or minutes off, so the app does not compare those
instants with the phone's clock directly. \code{my_status} also returns the server's clock, and each time
it is called (at launch and at every catch-up) the app records how far the phone's clock is off, taking
the server's answer as halfway through the round trip (\code{src/lib/server-clock.ts}).
\code{serverNow()}, the phone's clock corrected by that offset, drives the shared countdown clock, the
end-of-session timer, the wait for an answer, the query for missed requests and the earliest start of a
new event. Before the first answer the offset is zero, which was the behaviour before. Whatever the
phone's clock says, it cannot extend a session or a request, because the server checks its own time and
its minute job removes whatever has run out.

The countdowns themselves share one clock (\code{useClock}): a single one-second interval serves every
countdown on screen and stops when none is shown. It reads the time outside React's render, as React
requires, and copes with the device clock being changed.

\section{Staying reachable without push notifications}

Yap does not send push notifications, so a request reaches a user only while the app is open. Two choices
follow from that. While the user is live, the app keeps the screen awake, and the \ui{Go live} tab says so
(\ui{Keep Yap open to get requests}). And every screen that a request or a chat opens is a normal route
(\code{/incoming}, \code{/chat}), opened by the \code{(app)} layout, so a push notification added later
can open the same routes.

\section{Error handling}

Errors are handled at the edge where they can be explained. The API layer turns every refusal, Auth error
and network failure into a sentence. Screens catch the errors of the actions they start and show them in
an alert or inline. Store actions catch the failures that are expected while offline and leave them to the
next catch-up. The result is that a user always sees a short explanation, never a stack trace or database
text, and an action that failed can simply be tried again.

\section{Timers and schedules}

\Cref{tab:timers} lists every timer in the system.

\begin{table}[ht]
  \centering
  \caption{Timers and schedules.}
  \label{tab:timers}
  \tablefont
  \begin{tabularx}{\textwidth}{@{}L{48mm}L{40mm}Y@{}}
    \toprule
    \textsf{\textbf{Timer}} & \textsf{\textbf{Period or delay}} & \textsf{\textbf{Owner}} \\
    \midrule
    Housekeeping job & every minute & \code{pg_cron}, \code{expire_stale} \\
    Sign-in record pruning & daily at 03:00 & \code{pg_cron}, \code{prune_audit_log} \\
    Request answer window & 5 minutes & \code{chat_requests.expires_at} \\
    Nearby refresh & every 20\,s, foreground only & \code{useNearbyPolling} \\
    Answer poll & every 7\,s & \code{waitForAnswer} \\
    Answer timer & end time plus 3\,s & \code{waitForAnswer} \\
    End of live session & at the server's end time & \code{(app)/_layout} \\
    Countdown clock & every second, while shown & \code{useClock} \\
    Location fix & up to 8\,s & \code{useForegroundLocation}, map \\
    Nearby prompt & 8\,s on screen & \code{NearbyToast} \\
    Code resend & 60\,s cooldown & \code{(auth)/verify} \\
    Signed photo URL & 15 minutes; reused until 5 minutes remain; refusals retried after 60\,s & \code{signPaths} \\
    \bottomrule
  \end{tabularx}
\end{table}
